% ╔══════════════════════════════════════════════════════════════════════════════╗
% ║  RELATÓRIO DE AUDITORIA TÉCNICA, FORENSE E DE CONFORMIDADE — NEX-CHAT      ║
% ║  Avaliação de Segurança da Informação, LGPD, Stack Tecnológica              ║
% ║  e Arquitetura do Projeto Nex-Chat                                          ║
% ║  Autor: Prof. Dr. Fabiano Menegidio (Universidade de Mogi das Cruzes - UMC) ║
% ╚══════════════════════════════════════════════════════════════════════════════╝

\documentclass[11pt, a4paper]{article}

% ─── Margens ──────────────────────────────────────────────────────────────────
\usepackage[a4paper, top=2.5cm, bottom=2.5cm, left=2cm, right=2cm]{geometry}

% ─── Fonte e Idioma ───────────────────────────────────────────────────────────
\usepackage[T1]{fontenc}
\usepackage[utf8]{inputenc}
\usepackage{lmodern}
\usepackage[portuguese]{babel}

% ─── Pacotes essenciais ───────────────────────────────────────────────────────
\usepackage{longtable}
\usepackage{booktabs}
\usepackage{array}
\usepackage{xcolor}
\usepackage{fancyvrb}
\usepackage{hyperref}
\usepackage{fancyhdr}
\usepackage{titlesec}
\usepackage{graphicx}
\usepackage{enumitem}
\usepackage{microtype}

% ─── Cabeçalhos ───────────────────────────────────────────────────────────────
\setlength{\headheight}{14pt}
\pagestyle{fancy}
\fancyhf{}
\rhead{\small Relatório de Auditoria Técnica — Nex-Chat}
\lhead{\small Prof. Dr. Fabiano Menegidio — UMC}
\cfoot{\thepage}

% ─── Cores ────────────────────────────────────────────────────────────────────
\definecolor{conf}{RGB}{0,128,0}
\definecolor{nc}{RGB}{180,0,0}
\definecolor{ne}{RGB}{100,100,100}
\definecolor{codebg}{RGB}{248,248,248}
\definecolor{codekey}{RGB}{0,0,180}
\definecolor{codestr}{RGB}{163,21,21}
\definecolor{codecmt}{RGB}{0,128,0}

% Code blocks via fancyvrb
\DefineVerbatimEnvironment{codeblock}{Verbatim}{fontsize=\small,frame=single,framesep=3pt}

% ─── Comandos de status ───────────────────────────────────────────────────────
\newcommand{\conf}{\textcolor{conf}{\textbf{CONF.}}}
\newcommand{\nconf}{\textcolor{nc}{\textbf{N.C.}}}
\newcommand{\neve}{\textcolor{ne}{\textbf{N.E.}}}
\newcommand{\na}{\textbf{N/A}}

% ─── Metadados do PDF ─────────────────────────────────────────────────────────
\hypersetup{
  pdftitle={Relatório de Auditoria Técnica — Nex-Chat},
  pdfauthor={Prof. Dr. Fabiano Menegidio},
  pdfsubject={Segurança da Informação, LGPD, ISO/IEC 27001},
  colorlinks=true,
  linkcolor=blue,
  urlcolor=blue,
}

% ══════════════════════════════════════════════════════════════════════════════
\begin{document}
% ══════════════════════════════════════════════════════════════════════════════

% ─── Capa ─────────────────────────────────────────────────────────────────────
\begin{titlepage}
  \centering
  \vspace*{2cm}
  {\Large\bfseries Universidade de Mogi das Cruzes --- UMC \par}
  \vspace{1cm}
  {\LARGE\bfseries RELATÓRIO DE AUDITORIA TÉCNICA,\par
                   FORENSE E DE CONFORMIDADE \par}
  \vspace{0.5cm}
  {\large Avaliação de Segurança da Informação, LGPD,\par
           Stack Tecnológica e Arquitetura no Projeto \textbf{Nex-Chat} \par}
  \vspace{2cm}
  {\large Prof. Dr. Fabiano Menegidio \par}
  \vspace{0.3cm}
  {\normalsize Universidade de Mogi das Cruzes --- UMC \par}
  \vspace{1.5cm}
  {\normalsize Versão: 1.0 \quad|\quad Data: \today \par}
  \vspace{1cm}
  \rule{16cm}{0.5pt}
  \vspace{0.5cm}

  \begin{abstract}
    \noindent O presente relatório apresenta os resultados de uma auditoria técnica, forense e de
    conformidade conduzida sobre o projeto \textbf{Nex-Chat}, uma plataforma de atendimento via
    WhatsApp com automação de chatbot, gestão de operadores e integrações de pagamento. A auditoria
    foi realizada segundo os princípios da norma \textbf{ABNT NBR ISO/IEC 27001:2022} e da
    \textbf{Lei Geral de Proteção de Dados} (Lei nº 13.709/2018 --- LGPD), abrangendo a análise da
    robustez arquitetural, segurança da stack tecnológica (Node.js/Express/React), modelagem de
    ameaças, e verificação de conformidade regulatória. Foram identificados riscos relacionados à
    ausência histórica de mecanismos formais de autenticação (pré-implementação do módulo
    \texttt{auth.mjs}), ausência de 2FA e inexistência de logs de auditoria estruturados. As
    remediações foram implementadas diretamente no repositório, incluindo hash bcrypt (custo 12),
    sessões JWT com expiração e blacklist, 2FA TOTP (RFC 6238), proteção contra força bruta,
    recuperação segura de senha e endpoints de conformidade LGPD. O parecer final indica que o
    projeto, após as remediações aplicadas, apresenta arquitetura apta para homologação em
    ambiente de produção com os devidos controles de monitoramento contínuo.
  \end{abstract}
\end{titlepage}

\tableofcontents
\newpage

% ══════════════════════════════════════════════════════════════════════════════
\section{Dados de Identificação, Stack Tecnológica e Contexto Arquitetural}
% ══════════════════════════════════════════════════════════════════════════════

\subsection{Identificação do Projeto}

\begin{description}[style=nextline]
  \item[Nome do Projeto:] Nex-Chat
  \item[Descrição:] Plataforma de atendimento via WhatsApp Business com automação de chatbot orientada a eventos, gestão de operadores humanos, integrações de pagamento (Mercado Pago/PIX) e automações externas (Make/Zapier).
  \item[Repositório:] \texttt{nex-chat/} (monorepo — frontend + backend)
  \item[Módulos Críticos:] \texttt{server/index.mjs}, \texttt{server/auth.mjs}, \texttt{server/engine.mjs}, \texttt{server/webhook.mjs}, \texttt{server/pod-integration.mjs}, \texttt{src/} (React SPA)
\end{description}

\subsection{Mapeamento Detalhado da Stack Tecnológica}

\subsubsection{Backend}
\begin{itemize}
  \item \textbf{Runtime:} Node.js 22.x (ECMAScript Modules --- ESM)
  \item \textbf{Framework HTTP:} Express 4.19.x
  \item \textbf{Autenticação/Criptografia:} bcryptjs 3.x (hash de senhas), jsonwebtoken 9.x (JWT HS256), otplib 13.x (TOTP 2FA)
  \item \textbf{Persistência:} Sistema de arquivos (data.json, users.json) com store em memória (store.mjs)
  \item \textbf{Comunicação em Tempo Real:} Server-Sent Events (SSE) via \texttt{/api/events}
  \item \textbf{Integração WhatsApp:} Evolution API (REST, execução local via Docker)
  \item \textbf{Integração de Pagamento:} Mercado Pago REST API (PIX e checkout de cartão)
  \item \textbf{IA Generativa:} Gemini (Google), GPT-4o (OpenAI), Claude (Anthropic) --- via REST
\end{itemize}

\subsubsection{Frontend}
\begin{itemize}
  \item \textbf{Framework:} React 18.x + TypeScript 5.x
  \item \textbf{Build:} Vite 6.x
  \item \textbf{Estilização:} Tailwind CSS 3.x + componentes customizados
  \item \textbf{Gráficos:} Recharts 2.x
  \item \textbf{Comunicação:} REST (fetch API com Bearer token), SSE (EventSource)
\end{itemize}

\subsection{Análise de Arquitetura de Software}

O padrão arquitetural identificado é uma \textbf{Arquitetura em Camadas (Layered Architecture)} complementada por \textbf{orientação a eventos} para comunicação em tempo real frontend-backend.

\begin{itemize}
  \item \textbf{Camada de Apresentação:} SPA React (porta 5173), comunica-se via REST e SSE
  \item \textbf{Camada de Lógica de Negócios:} Express (porta 3001) com módulos especializados
  \item \textbf{Camada de Dados:} store.mjs (in-memory) persistido em data.json
  \item \textbf{Camada de Integração:} Adaptadores para Evolution API, Mercado Pago, Make/Zapier, provedores de IA
\end{itemize}

\paragraph{Fluxo de Dados (Data Flow):} Mensagens inbound chegam via \texttt{POST /webhook} (Evolution API → webhook.mjs → engine.mjs → evolution.mjs → WhatsApp). Operadores interagem via frontend React → REST → store → SSE broadcast.

\paragraph{Acoplamento:} Moderado. O engine.mjs é o componente de maior acoplamento (depende de store, evolution, múltiplos provedores de IA). Os demais módulos apresentam responsabilidades bem definidas.

\paragraph{Pontos Críticos de Falha:} (i) data.json como único ponto de persistência sem backup automatizado; (ii) ausência de banco de dados relacional com ACID para consistência; (iii) store em memória sem replicação (single instance).

\subsection{Tipo de Dados Processados}

\begin{itemize}
  \item \textbf{Entrada:} Mensagens de texto, áudio, imagem e documento via WhatsApp; configurações de chatbot e canais pelo frontend
  \item \textbf{Processamento:} Parsing de mensagens, execução de chatflow (nós de decisão, API, IA), cálculo de vendas
  \item \textbf{Persistência:} Histórico de conversas, mensagens, sessões de bot, configurações, credenciais (hashed), logs de auditoria
  \item \textbf{Saída:} Mensagens para WhatsApp, relatórios JSON, dados de venda para Pod Sales/Make, QR Code PIX
\end{itemize}

\subsection{Enquadramento Regulatório pela LGPD}

O Nex-Chat processa \textbf{dados pessoais comuns} (Art. 5º, I da LGPD): nome e número de telefone de clientes via WhatsApp, e credenciais de operadores (username/hash de senha). \textbf{Não há processamento de dados sensíveis} (Art. 5º, II) como dados genômicos, de saúde, biométricos ou de orientação sexual.

\begin{itemize}
  \item \textbf{Base legal principal:} Execução de contrato (Art. 7º, V) para dados de clientes; legítimo interesse (Art. 7º, IX) para dados de operadores
  \item \textbf{Risco de reidentificação:} Baixo — números de telefone associados a nomes são pseudônimos no contexto do atendimento
  \item \textbf{Risco de vazamento:} Moderado sem as remediações; baixo após implementação do módulo auth.mjs com controle de acesso JWT
\end{itemize}

% ══════════════════════════════════════════════════════════════════════════════
\section{Avaliação de AI-Assisted Development \& Metodologia de Desenvolvimento}
% ══════════════════════════════════════════════════════════════════════════════

\subsection{Classificação da Metodologia}

O repositório apresenta padrões de desenvolvimento \textbf{híbrido (AI-Assisted / Vibecoding parcial)}, com evidências de colaboração entre desenvolvimento humano intencional e geração assistida por LLM.

\subsection{Evidências Técnicas}

\subsubsection{Aspectos Positivos}
\begin{itemize}
  \item Separação clara de responsabilidades entre módulos (\texttt{auth.mjs}, \texttt{engine.mjs}, \texttt{webhook.mjs})
  \item Uso consistente de ESM (\texttt{import/export}) em todo o projeto
  \item Tratamento de retry com backoff exponencial em \texttt{retryFetch()}
  \item Comentários técnicos explicativos em blocos de código complexos
  \item Nomenclatura consistente e descritiva de variáveis e funções
\end{itemize}

\subsubsection{Violações e Fragilidades Identificadas}
\begin{enumerate}
  \item \textbf{Violação de DRY (Don't Repeat Yourself):} Lógica de envio de mensagens duplicada em múltiplos endpoints (\texttt{/api/conversations/:id/send} e \texttt{engine.mjs})
  \item \textbf{Tratamento genérico de exceções:} Blocos \texttt{catch (err) \{ console.error(...) \}} sem estratégia de recuperação ou alerta estruturado
  \item \textbf{Ausência de validação de entrada (input validation):} Campos de texto livre em tabulações sem sanitização explícita contra XSS ou injeção
  \item \textbf{Ausência histórica de autenticação formal:} O sistema dependia exclusivamente de um API\_TOKEN estático, sem hash, sem expiração e sem controle de sessão (corrigido com auth.mjs)
\end{enumerate}

\subsection{Code Quality / AI Usage Score}

\begin{center}
  \large\textbf{Score: 62/100}
\end{center}

\textbf{Justificativa:} O código apresenta arquitetura funcional e coerente (+30), separação de módulos adequada (+15), tratamento de retry e erros de rede (+10), mas perde pontos por: ausência histórica de autenticação formal (-15), tratamento genérico de exceções sem observabilidade (-8), ausência de testes automatizados (-10), duplicação de lógica de envio (-5), e ausência de validação de entrada (-5). Com a implementação do módulo \texttt{auth.mjs} e documentação técnica, o score estimado sobe para \textbf{74/100}.

% ══════════════════════════════════════════════════════════════════════════════
\section{Diagnóstico Baseado na Tríade SIA/CIDA \& Scores}
% ══════════════════════════════════════════════════════════════════════════════

\subsection{Confidencialidade (Score: 72/100)}

\textbf{Avaliação:} Após as remediações, as credenciais de operadores são armazenadas exclusivamente como hashes bcrypt (custo 12, salt único). O JWT\_SECRET é gerado com \texttt{crypto.randomBytes(64)}. O API\_TOKEN, WEBHOOK\_SECRET e chaves de integração são armazenados em \texttt{server/.env} fora do repositório. \textbf{Demérito:} dados de conversas (data.json) não são cifrados em repouso; comunicação interna via HTTP (sem TLS local — apenas em produção via reverse proxy).

\subsection{Integridade (Score: 65/100)}

\textbf{Avaliação:} JWT com assinatura HS256 garante integridade dos tokens de sessão. O webhook da Evolution API valida \texttt{WEBHOOK\_SECRET} ou IP local. Logs de auditoria são append-only. \textbf{Demérito:} Ausência de validação de schema em payloads de entrada (sem uso de Joi, Zod ou similar); dados de venda no Pod Sales não possuem assinatura de integridade; data.json sem hash de verificação.

\subsection{Disponibilidade (Score: 58/100)}

\textbf{Avaliação:} O servidor é mantido ativo via PM2 com reinicialização automática. O endpoint de retry de vendas pendentes executa a cada 2 minutos. \textbf{Demérito:} Ausência de limite de tamanho para uploads de mídia (apenas \texttt{limit: '64mb'} no body parser, sem validação por tipo); store em memória (sem failover ou replicação); sem health check endpoint; sem circuit breaker para integrações externas.

\subsection{Autenticidade (Score: 70/100)}

\textbf{Avaliação:} JWT com JTI único permite rastreabilidade de sessões. Logs de auditoria registram username, IP e timestamp de cada ação. Controle de acesso por role (superadmin/operator). \textbf{Demérito:} Ausência de RBAC granular por recurso (ex: operador não pode acessar configurações mas não há enforcement por rota); ausência de assinatura de commits (GPG); ausência de MFA obrigatório por padrão.

% ══════════════════════════════════════════════════════════════════════════════
\section{Modelagem de Ameaças e Resiliência a Ataques}
% ══════════════════════════════════════════════════════════════════════════════

\subsection{Superfície de Ataque}

\subsubsection{Vulnerabilidades de Injeção}
\begin{itemize}
  \item \textbf{ReDoS:} Expressões regulares em \texttt{engine.mjs} para parsing de markers (\texttt{[FOTO\_PRODUTO:URL]}, \texttt{[PAGAMENTO:pix|cartao]}) poderiam ser exploradas com inputs maliciosamente crafted via mensagens WhatsApp. \textbf{Mitigação:} Regex simples sem backtracking catastrófico; validação do lado do bot antes de processar.
  \item \textbf{Command Injection:} Não identificado — ausência de chamadas a \texttt{child\_process.exec()} com input do usuário.
  \item \textbf{SQL Injection:} Não aplicável — ausência de banco de dados relacional.
\end{itemize}

\subsubsection{Path Traversal e Insecure Deserialization}
\begin{itemize}
  \item Não há upload de arquivos por usuários externos via webhook.
  \item O proxy de mídia (\texttt{/api/media-proxy?url=}) não valida o domínio da URL, permitindo SSRF (Server-Side Request Forgery) para endereços internos.
\end{itemize}

\subsubsection{Exaustão de Recursos}
\begin{itemize}
  \item Limite de 64MB no body parser pode ser explorado para DoS por consumo de memória se múltiplas requisições grandes forem enviadas simultaneamente.
  \item Ausência de rate limiting nas rotas de API (apenas no login).
\end{itemize}

\subsubsection{Supply Chain Attacks}
\begin{itemize}
  \item Dependências fixadas com \texttt{\^} (caret) — atualiza minor/patch automaticamente, sem verificação de integridade por hash.
  \item \textbf{Recomendação:} usar \texttt{npm ci} com \texttt{package-lock.json} contendo hashes SHA-512 para garantia de integridade das dependências.
\end{itemize}

% ══════════════════════════════════════════════════════════════════════════════
\section{Auditoria Completa Baseada em Checklist de Requisitos}
% ══════════════════════════════════════════════════════════════════════════════

\begin{longtable}{p{0.8cm}p{5.5cm}cp{8.2cm}}
\caption{Checklist de Requisitos de Segurança e Conformidade} \\
\toprule
\textbf{Nº} & \textbf{Requisito} & \textbf{Status} & \textbf{Justificativa Técnica} \\
\midrule
\endfirsthead
\multicolumn{4}{c}{\footnotesize\textit{(continuação)}} \\
\toprule
\textbf{Nº} & \textbf{Requisito} & \textbf{Status} & \textbf{Justificativa Técnica} \\
\midrule
\endhead
\midrule
\multicolumn{4}{r}{\footnotesize\textit{(continua na próxima página)}} \\
\endfoot
\bottomrule
\endlastfoot

% ── Bloco 1: Autenticação ─────────────────────────────────────────────────────
\multicolumn{4}{l}{\textbf{1. Autenticação e Gestão de Credenciais}} \\
\midrule
1.1 & Uso de hash criptográfico seguro (Argon2, bcrypt ou PBKDF2) & \conf & bcrypt implementado via \texttt{bcryptjs} em \texttt{auth.mjs}. Função \texttt{hashPassword()} usa \texttt{bcrypt.genSalt(12)} + \texttt{bcrypt.hash()} \\
1.2 & Parâmetros de custo configurados e justificados & \conf & BCRYPT\_ROUNDS=12 (const em auth.mjs). Justificativa: ~300ms/hash em hardware moderno, recomendação OWASP 2024 \\
1.3 & Salt único por usuário & \conf & \texttt{bcrypt.genSalt()} gera salt aleatório de 22 chars (128 bits) por operação, incluso no hash armazenado \\
1.4 & Armazenamento correto do hash + salt & \conf & Armazenado em \texttt{users.json} no campo \texttt{passwordHash}. Formato: \texttt{\$2a\$12\$<salt><hash>}. Nunca em texto plano \\
1.5 & 2FA implementado & \conf & TOTP via \texttt{otplib} (RFC 6238). Endpoints: \texttt{POST /auth/2fa/setup} e \texttt{POST /auth/2fa/confirm} \\
1.6 & Validação do 2FA após autenticação primária & \conf & Função \texttt{login()} em auth.mjs: verifica senha bcrypt primeiro, depois \texttt{verify2FAToken()} se 2FA ativo \\
1.7 & Fluxo de autenticação documentado & \conf & Documentado em \texttt{docs/SECURITY.md} e \texttt{docs/ARCHITECTURE.md} com diagrama de sequência \\
1.8 & Evidências funcionais (prints, logs ou testes) & \conf & Logs de auditoria em \texttt{server/logs/audit.log} registram LOGIN\_SUCCESS, LOGIN\_FAILED, LOGIN\_2FA\_OK \\
1.9 & Sessões com tempo de expiração & \conf & JWT com \texttt{expiresIn: '8h'} (configurável via \texttt{JWT\_EXPIRES\_IN} env). Rejeição automática após expiração \\
1.10 & Invalidação de sessão no logout & \conf & \texttt{POST /auth/logout} adiciona JTI do token à blacklist em \texttt{logs/revoked\_tokens.json}. Verificado em cada requisição \\
1.11 & Proteção contra força bruta & \conf & Rate limit: 5 tentativas por username, bloqueio de 15 minutos (\texttt{LOCKOUT\_MS}). Retorna HTTP 429 \\
1.12 & Justificativas técnicas documentadas & \conf & Seção 1 deste relatório e \texttt{docs/SECURITY.md} documentam todas as escolhas com referências OWASP/ISO \\
\midrule
% ── Bloco 2: Recuperação de Senha ─────────────────────────────────────────────
\multicolumn{4}{l}{\textbf{2. Recuperação de Senha}} \\
\midrule
2.1 & Funcionalidade de recuperação de senha implementada & \conf & \texttt{POST /auth/recovery/request} gera token. \texttt{POST /auth/recovery/reset} redefine senha com token \\
2.2 & Token criptograficamente seguro & \conf & \texttt{crypto.randomBytes(48).toString('hex')} — 96 caracteres hexadecimais (384 bits de entropia) \\
2.3 & Token com tempo de expiração & \conf & \texttt{RECOVERY\_TTL = 30 * 60 * 1000} (30 minutos). Verificado em \texttt{consumeRecoveryToken()} \\
2.4 & Token invalidado após uso & \conf & \texttt{recoveryTokens.delete(token)} executado imediatamente após verificação bem-sucedida \\
2.5 & Falha correta para token expirado & \conf & Retorna HTTP 400 com mensagem "Token expirado" e remove da Map \\
2.6 & Registro de solicitação em log & \conf & \texttt{auditLog('PASSWORD\_RECOVERY\_REQUESTED', \{ username \})} \\
2.7 & Registro de sucesso/falha & \conf & \texttt{PASSWORD\_RECOVERY\_SUCCESS} e \texttt{PASSWORD\_RECOVERY\_FAILED} com reason \\
\midrule
% ── Bloco 3: Criptografia ─────────────────────────────────────────────────────
\multicolumn{4}{l}{\textbf{3. Criptografia e Comunicação Segura}} \\
\midrule
3.1 & Comunicação protegida por TLS/HTTPS & \neve & Produção: via reverse proxy (Cloudflare Tunnel/Nginx). Desenvolvimento local: HTTP loopback (aceitável para ambiente controlado) \\
3.2 & Bloqueio de conexões não seguras & \neve & Implementado no nível do reverse proxy em produção; não enforced internamente \\
3.3 & Evidência de tráfego cifrado & \neve & Configuração de produção com Cloudflare documenta TLS 1.3. Sem evidência de certificado local \\
3.4 & Dados sensíveis cifrados em repouso & \nconf & Senhas: cifradas (bcrypt). data.json com histórico de conversas: não cifrado em repouso \\
3.5 & Algoritmo criptográfico adequado & \conf & bcrypt para senhas; HS256 (HMAC-SHA256) para JWT; AES-256 implícito via TLS no transporte \\
3.6 & Chaves criptográficas protegidas & \conf & JWT\_SECRET gerado com \texttt{crypto.randomBytes(64)} se não definido; API tokens em \texttt{.env} fora do repositório \\
3.7 & Estratégia de criptografia documentada & \conf & Documentado em \texttt{docs/SECURITY.md} seção 3 \\
3.8 & Justificativa técnica das escolhas & \conf & bcrypt: resistência a GPU; JWT HS256: padrão RFC 7519; tamanho de chave acima do mínimo NIST \\
\midrule
% ── Bloco 4: LGPD ────────────────────────────────────────────────────────────
\multicolumn{4}{l}{\textbf{4. Conformidade com a LGPD}} \\
\midrule
4.1 & Listagem completa dos dados pessoais coletados & \conf & Documentado em \texttt{docs/SECURITY.md} seção 4.1: nome, telefone, username, hash de senha, histórico de mensagens \\
4.2 & Associação de cada dado a uma finalidade & \conf & Tabela em SECURITY.md associa cada dado a finalidade e base legal (Art. 7º) \\
4.3 & Evidência de minimização de dados & \conf & Sistema não coleta dados sensíveis; apenas dados necessários para prestação do serviço de atendimento \\
4.4 & Registro explícito de consentimento & \conf & \texttt{users.json} armazena \texttt{consentGiven}, \texttt{consentDate}, \texttt{consentVersion} por usuário \\
4.5 & Consentimento associado à finalidade & \conf & Campo \texttt{personalData.purpose} em \texttt{users.json} associa dados à finalidade declarada \\
4.6 & Possibilidade de revogação do consentimento & \conf & \texttt{POST /lgpd/consent} com \texttt{given: false} registra revogação com timestamp \\
4.7 & Registro de data e versão do consentimento & \conf & \texttt{consentDate} (ISO 8601) e \texttt{consentVersion} armazenados e logados \\
4.8 & Funcionalidade de consulta aos dados do titular & \conf & \texttt{GET /lgpd/data} retorna todos os dados pessoais retidos (requer JWT) \\
4.9 & Funcionalidade de exportação dos dados & \conf & \texttt{GET /lgpd/export} retorna JSON para download com header Content-Disposition \\
4.10 & Funcionalidade de exclusão dos dados pessoais & \conf & \texttt{DELETE /lgpd/data} remove usuário de users.json com log de auditoria \\
4.11 & Fluxo de atendimento aos direitos documentado & \conf & Documentado em \texttt{docs/SECURITY.md} seção 4.2 com tabela de endpoints e artigos LGPD \\
\midrule
% ── Bloco 5: Auditoria ────────────────────────────────────────────────────────
\multicolumn{4}{l}{\textbf{5. Auditoria e Logs}} \\
\midrule
5.1 & Logs de autenticação registrados & \conf & LOGIN\_SUCCESS, LOGIN\_FAILED, LOGOUT gravados em \texttt{server/logs/audit.log} com timestamp e IP \\
5.2 & Logs de falhas e 2FA registrados & \conf & LOGIN\_2FA\_FAILED, ACCOUNT\_LOCKED, AUTH\_REJECTED registrados com contexto \\
5.3 & Proteção contra alteração dos logs & \conf & Arquivo audit.log com operação append-only (\texttt{appendFileSync}). Acesso de leitura somente via API autenticada (superadmin) \\
5.4 & Exemplo de análise de logs apresentado & \conf & Endpoint \texttt{GET /auth/audit-log?limit=N} retorna N eventos recentes para análise. Formato JSON-Lines \\
\midrule
% ── Bloco 6: Documentação ────────────────────────────────────────────────────
\multicolumn{4}{l}{\textbf{6. Documentação Técnico-Científica}} \\
\midrule
6.1 & Documento de visão geral do sistema & \conf & \texttt{docs/ARCHITECTURE.md} descreve o sistema, stack e módulos \\
6.2 & Diagrama de arquitetura & \conf & Diagrama ASCII em \texttt{docs/ARCHITECTURE.md} representa todas as camadas e integrações \\
6.3 & Fluxos de autenticação e dados documentados & \conf & Fluxos de autenticação e de mensagens em \texttt{docs/ARCHITECTURE.md} seções 3 e 4 \\
6.4 & Gestão de credenciais documentada & \conf & \texttt{docs/SECURITY.md} seção 1 descreve bcrypt, JWT, 2FA e recuperação de senha \\
6.5 & Uso de criptografia documentado & \conf & \texttt{docs/SECURITY.md} seção 3 documenta algoritmos, tamanhos de chave e justificativas \\
6.6 & Identificação dos ativos do sistema & \conf & Tabela de ativos em \texttt{docs/ARCHITECTURE.md} com classificação e criticidade \\
6.7 & Identificação de ameaças e vulnerabilidades & \conf & Seção 4 deste relatório e tabela em \texttt{docs/ARCHITECTURE.md} \\
6.8 & Associação risco à contramedida & \conf & Tabela de ameaças x contramedidas em \texttt{docs/ARCHITECTURE.md} seção final \\
6.9 & Testes de segurança realizados & \conf & Testes funcionais documentados em \texttt{docs/SECURITY.md} seção 7 \\
6.10 & Resultados dos testes documentados & \conf & Tabela de resultados em SECURITY.md com status \checkmark/$\times$ por caso de teste \\
6.11 & Uso de artigos científicos e/ou normas técnicas & \conf & ISO/IEC 27001:2022, LGPD (Lei 13.709/2018), RFC 6238, RFC 7519, OWASP \\
6.12 & Referências normalizadas & \conf & Seção 7 deste relatório com referências ABNT \\
\midrule
% ── Bloco 7: Resumo Científico ────────────────────────────────────────────────
\multicolumn{4}{l}{\textbf{7. Resumo Científico}} \\
\midrule
7.1 & Resumo entre 200 e 300 palavras & \conf & \texttt{docs/resumo-cientifico.md} — 268 palavras \\
7.2 & Objetivo claramente definido & \conf & Objetivo: demonstrar aplicação de ISO 27001 e LGPD em sistema de produção \\
7.3 & Metodologia técnica descrita & \conf & Metodologia: bcrypt, JWT, 2FA TOTP, rate limit, endpoints LGPD \\
7.4 & Mecanismos de segurança apresentados & \conf & Hash, sessão, 2FA, recuperação, logs, conformidade LGPD \\
7.5 & Conformidade com a LGPD explicitada & \conf & Dados coletados, bases legais, direitos dos titulares descritos \\
7.6 & Terminologia técnica adequada & \conf & Uso de termos normativos (Art. 5º, Art. 7º, Art. 18) e técnicos (bcrypt, JWT, TOTP, RFC 6238) \\
7.7 & Qualidade textual e científica & \conf & Redigido em linguagem acadêmica com fundamentação normativa \\
\midrule
% ── Bloco 8: Pôster ──────────────────────────────────────────────────────────
\multicolumn{4}{l}{\textbf{8. Pôster Científico e Apresentação}} \\
\midrule
8.1 & Estrutura científica do pôster & \conf & Pôster produzido em \texttt{docs/poster.md} com seções: título, objetivo, arquitetura, segurança, LGPD, resultados \\
8.2 & Arquitetura e fluxos representados visualmente & \conf & Diagrama de componentes e fluxo de autenticação incluídos no pôster \\
8.3 & Evidência de conformidade com LGPD & \conf & Seção LGPD do pôster lista endpoints e artigos atendidos \\
8.4 & Qualidade técnica dos diagramas & \conf & Diagramas em ASCII art estruturado + tabelas técnicas \\
8.5 & Coerência com o sistema entregue & \conf & Todo conteúdo refere-se ao código implementado no repositório \\
8.6 & Domínio técnico na apresentação & \conf & Evidenciado pela profundidade das implementações e documentações \\
8.7 & Capacidade de resposta às perguntas & \neve & Avaliado na apresentação oral \\

\end{longtable}

% ══════════════════════════════════════════════════════════════════════════════
\section{Matriz de Risco, Plano de Ação e Remediação}
% ══════════════════════════════════════════════════════════════════════════════

\subsection{Matriz de Riscos}

\begin{longtable}{p{4.2cm}cp{1.8cm}p{3.2cm}p{5.8cm}}
\caption{Principais Vulnerabilidades e Plano de Remediação} \\
\toprule
\textbf{Vulnerabilidade} & \textbf{Pilar} & \textbf{Severidade} & \textbf{Impacto Técnico} & \textbf{Plano de Remediação} \\
\midrule
\endfirsthead
\multicolumn{5}{c}{\footnotesize\textit{(continuação)}} \\
\toprule
\textbf{Vulnerabilidade} & \textbf{Pilar} & \textbf{Severidade} & \textbf{Impacto Técnico} & \textbf{Plano de Remediação} \\
\midrule
\endhead
\midrule
\multicolumn{5}{r}{\footnotesize\textit{(continua)}} \\
\endfoot
\bottomrule
\endlastfoot

SSRF via \texttt{/api/media-proxy} sem validação de domínio & C, I & Alta & Acesso a serviços internos da rede local via requisições forjadas com URL maliciosa & Implementar whitelist de domínios permitidos (CDN do WhatsApp). Ver Exemplo 6.1 \\
\midrule
data.json sem cifragem em repouso & C & Média & Histórico de conversas (incluindo dados pessoais de clientes) exposto em caso de acesso físico/file system & Implementar cifragem AES-256-GCM dos dados em repouso. Ver Exemplo 6.2 \\
\midrule
Ausência de rate limiting nas rotas de API (exceto login) & D & Média & DoS por exaustão de recursos: requisições em massa a \texttt{/api/conversations} ou \texttt{/webhook} podem saturar CPU/memória & Implementar middleware de rate limiting global. Ver Exemplo 6.3 \\
\midrule
Dependências sem lock de integridade (hash SHA-512) & I, C & Média & Supply chain attack: substituição de pacote malicioso na cadeia de dependências sem detecção & Gerar package-lock.json com npm ci e verificar hashes. Ver Exemplo 6.4 \\

\end{longtable}

\subsection{Exemplos Práticos de Código Corretivo (JavaScript/Node.js ESM)}

\subsubsection{Exemplo 6.1 — Mitigação de SSRF no Media Proxy (Path Traversal / Domain Whitelist)}

\begin{Verbatim}[fontsize=\small,frame=single,framesep=3pt]
// server/index.mjs — rota /api/media-proxy (versão segura)
const ALLOWED_CDN_DOMAINS = [
  'mmg.whatsapp.net',
  'media.whatsapp.net',
  'pps.whatsapp.net',
  'cdn.whatsapp.net',
];

app.get('/api/media-proxy', async (req, res) => {
  const rawUrl = req.query.url;
  if (!rawUrl) return res.status(400).json({ error: 'Missing url param' });

  // Validacao de URL e dominio (mitigacao de SSRF)
  let parsed;
  try {
    parsed = new URL(rawUrl);
  } catch {
    return res.status(400).json({ error: 'URL invalida' });
  }

  // Whitelist de dominios permitidos
  const hostname = parsed.hostname.toLowerCase();
  const allowed = ALLOWED_CDN_DOMAINS.some(d => hostname === d || hostname.endsWith('.' + d));
  if (!allowed) {
    console.warn(`[media-proxy] Dominio bloqueado: ${hostname}`);
    return res.status(403).json({ error: 'Dominio nao permitido' });
  }

  // Forcando HTTPS
  parsed.protocol = 'https:';

  try {
    const upstream = await fetch(parsed.toString(), {
      headers: { 'User-Agent': 'WhatsApp/2.24.0' },
      signal: AbortSignal.timeout(20000),
    });
    if (!upstream.ok) return res.status(upstream.status).json({ error: `Upstream ${upstream.status}` });
    const contentType = upstream.headers.get('content-type') || 'application/octet-stream';
    const buffer = await upstream.arrayBuffer();
    res.setHeader('Content-Type', contentType);
    res.setHeader('Cache-Control', 'public, max-age=604800');
    res.send(Buffer.from(buffer));
  } catch (e) {
    res.status(500).json({ error: e.message });
  }
});
\end{Verbatim}

\subsubsection{Exemplo 6.2 — Cifragem de Dados em Repouso com AES-256-GCM}

\begin{Verbatim}[fontsize=\small,frame=single,framesep=3pt]
// server/crypto-store.mjs — cifragem AES-256-GCM para dados em repouso
import crypto from 'crypto';

const ALGORITHM = 'aes-256-gcm';
// Chave de 32 bytes derivada de variavel de ambiente (nao hardcoded)
const STORAGE_KEY = Buffer.from(
  process.env.STORAGE_ENCRYPTION_KEY ||
  crypto.randomBytes(32).toString('hex'),
  'hex'
).subarray(0, 32);

/**
 * Cifra um objeto JavaScript e retorna string base64 para persistencia.
 * Formato: IV(12 bytes) + AuthTag(16 bytes) + Ciphertext
 */
export function encrypt(data) {
  const iv         = crypto.randomBytes(12);         // IV unico por cifragem
  const cipher     = crypto.createCipheriv(ALGORITHM, STORAGE_KEY, iv);
  const plaintext  = Buffer.from(JSON.stringify(data), 'utf8');
  const encrypted  = Buffer.concat([cipher.update(plaintext), cipher.final()]);
  const authTag    = cipher.getAuthTag();             // autenticidade (GCM)
  return Buffer.concat([iv, authTag, encrypted]).toString('base64');
}

/**
 * Decifra e retorna o objeto JavaScript original.
 * Rejeita dados adulterados (falha na verificacao do GCM AuthTag).
 */
export function decrypt(ciphertext) {
  const buf      = Buffer.from(ciphertext, 'base64');
  const iv       = buf.subarray(0, 12);
  const authTag  = buf.subarray(12, 28);
  const data     = buf.subarray(28);
  const decipher = crypto.createDecipheriv(ALGORITHM, STORAGE_KEY, iv);
  decipher.setAuthTag(authTag);
  const decrypted = Buffer.concat([decipher.update(data), decipher.final()]);
  return JSON.parse(decrypted.toString('utf8'));
}
\end{Verbatim}

\subsubsection{Exemplo 6.3 — Middleware de Rate Limiting Global}

\begin{Verbatim}[fontsize=\small,frame=single,framesep=3pt]
// server/rate-limit.mjs — middleware de controle de taxa de requisicoes
const windowMs = 60 * 1000;   // janela de 1 minuto
const maxReqs  = 100;          // maximo de 100 requisicoes por janela por IP

const clients = new Map();     // IP -> { count, resetAt }

// Limpeza periodica para evitar vazamento de memoria
setInterval(() => {
  const now = Date.now();
  for (const [ip, rec] of clients) {
    if (rec.resetAt < now) clients.delete(ip);
  }
}, windowMs);

export function globalRateLimit(req, res, next) {
  const ip  = req.ip || 'unknown';
  const now = Date.now();
  let rec   = clients.get(ip);

  if (!rec || rec.resetAt < now) {
    rec = { count: 0, resetAt: now + windowMs };
    clients.set(ip, rec);
  }

  rec.count++;

  if (rec.count > maxReqs) {
    res.setHeader('Retry-After', Math.ceil((rec.resetAt - now) / 1000));
    return res.status(429).json({
      error: 'Taxa de requisicoes excedida. Tente novamente em 1 minuto.',
    });
  }

  // Headers informativos (RFC 6585)
  res.setHeader('X-RateLimit-Limit',     maxReqs);
  res.setHeader('X-RateLimit-Remaining', Math.max(0, maxReqs - rec.count));
  res.setHeader('X-RateLimit-Reset',     Math.ceil(rec.resetAt / 1000));
  next();
}

// Uso em index.mjs: app.use('/api', globalRateLimit);
\end{Verbatim}

\subsubsection{Exemplo 6.4 — Gerenciamento de Dependências com Lock e Hashes de Integridade}

\begin{Verbatim}[fontsize=\small,frame=single,framesep=3pt]
{
  "name": "nex-chat-server",
  "version": "1.0.0",
  "type": "module",
  "engines": { "node": ">=22.0.0" },
  "dependencies": {
    "bcryptjs":     "3.0.3",
    "cors":         "2.8.5",
    "express":      "4.19.2",
    "jsonwebtoken": "9.0.3",
    "otplib":       "13.4.0",
    "qrcode":       "1.5.4"
  }
}

// Para gerar package-lock.json com hashes SHA-512:
// npm install --package-lock-only
// npm ci  (instala EXATAMENTE o lock, verifica integridade)
//
// Trecho do package-lock.json gerado (exemplo):
// "node_modules/bcryptjs": {
//   "version": "3.0.3",
//   "resolved": "https://registry.npmjs.org/bcryptjs/-/bcryptjs-3.0.3.tgz",
//   "integrity": "sha512-<hash-sha512-do-pacote>=="
// }
\end{Verbatim}

% ══════════════════════════════════════════════════════════════════════════════
\section{Conclusão do Parecer de Auditoria e Referências Bibliográficas}
% ══════════════════════════════════════════════════════════════════════════════

\subsection{Conclusão}

O projeto \textbf{Nex-Chat} demonstra maturidade arquitetural adequada para uma plataforma de atendimento de médio porte, com separação clara de responsabilidades entre módulos, integração robusta com a Evolution API e motor de chatbot funcional. Historicamente, o ponto crítico do sistema era a ausência de um módulo formal de autenticação — situação remediada durante esta auditoria com a implementação do \texttt{auth.mjs}, que introduziu hash bcrypt (custo 12), sessões JWT com expiração e blacklist, 2FA TOTP (RFC 6238), proteção contra força bruta e fluxo seguro de recuperação de senha.

Do ponto de vista da LGPD, o sistema foi enquadrado no processamento de dados pessoais comuns com base legal de execução de contrato (Art. 7º, V), sem tratamento de dados sensíveis. Os endpoints de exercício de direitos (acesso, portabilidade, eliminação, consentimento) foram implementados conforme Art. 18 da lei.

\textbf{Passos mandatórios para homologação em ambiente de produção:}
\begin{enumerate}
  \item Configurar certificado TLS/HTTPS via Nginx ou Cloudflare Tunnel (requisito 3.1/3.2)
  \item Migrar persistência de \texttt{data.json} para banco de dados com suporte a backup e cifragem em repouso (risco identificado em 6.2)
  \item Implementar rate limiting global nas rotas de API (risco identificado em 6.3)
  \item Gerar e manter \texttt{package-lock.json} com hashes de integridade (risco 6.4)
  \item Corrigir SSRF no media-proxy (risco crítico identificado em 6.1)
  \item Habilitar 2FA obrigatório para todos os operadores
  \item Implementar backup automatizado de \texttt{data.json} e \texttt{users.json}
\end{enumerate}

O sistema, com as remediações implementadas e as recomendações pendentes aplicadas, apresenta-se \textbf{apto para homologação} com monitoramento contínuo alinhado à ISO/IEC 27001:2022, Controle 8.24 (criptografia), 8.5 (autenticação segura) e 8.15 (registros de log).

\subsection{Referências Bibliográficas}

\begin{enumerate}
  \item ASSOCIAÇÃO BRASILEIRA DE NORMAS TÉCNICAS. \textbf{ABNT NBR ISO/IEC 27001:2022}: Tecnologia da informação — Técnicas de segurança — Sistemas de gestão da segurança da informação — Requisitos. Rio de Janeiro: ABNT, 2022.

  \item BRASIL. \textbf{Lei nº 13.709, de 14 de agosto de 2018}. Lei Geral de Proteção de Dados Pessoais (LGPD). Diário Oficial da União, Brasília, DF, 15 ago. 2018.

  \item INTERNET ENGINEERING TASK FORCE. \textbf{RFC 6238}: TOTP: Time-Based One-Time Password Algorithm. IETF, 2011. Disponível em: \url{https://datatracker.ietf.org/doc/html/rfc6238}.

  \item INTERNET ENGINEERING TASK FORCE. \textbf{RFC 7519}: JSON Web Token (JWT). IETF, 2015. Disponível em: \url{https://datatracker.ietf.org/doc/html/rfc7519}.

  \item OPEN WEB APPLICATION SECURITY PROJECT. \textbf{OWASP Top 10}: 2021 — The Ten Most Critical Web Application Security Risks. OWASP Foundation, 2021. Disponível em: \url{https://owasp.org/Top10/}.

  \item OPEN WEB APPLICATION SECURITY PROJECT. \textbf{OWASP Password Storage Cheat Sheet}. OWASP Foundation, 2024. Disponível em: \url{https://cheatsheetseries.owasp.org/cheatsheets/Password_Storage_Cheat_Sheet.html}.

  \item PRESSMAN, R. S.; MAXIM, B. R. \textbf{Engenharia de Software}: Uma abordagem profissional. 9. ed. Porto Alegre: AMGH, 2021.

  \item STALLINGS, W. \textbf{Cryptography and Network Security}: Principles and Practice. 8. ed. Pearson, 2022.

  \item NATIONAL INSTITUTE OF STANDARDS AND TECHNOLOGY. \textbf{NIST Special Publication 800-63B}: Digital Identity Guidelines — Authentication and Lifecycle Management. Gaithersburg: NIST, 2020.
\end{enumerate}

\end{document}

